\chapter{太阳–地球–月球系统}\label{ch:g6:sun-earth-moon}

主宰我们天空的三个天体：给我们白昼的太阳、
承载我们的地球、陪伴我们的\omterm{def:g2:moon-in-the-sky:moon}{月球}。你对它们都很熟悉；
现在我们来勾勒这个\emph{系统}——谁绕谁转、
转得多快、相距多远——并还清一笔旧账：
\omterm{def:g3:sun-days-seasons:year}{季节}的秘密，就在本章之中。

\section{三种运动，三只时钟}

\begin{definition}[地轴与轨道]\label{def:g6:sun-earth-moon:axisorbit}
地球的\emph{地轴}\index{地轴}是穿过两极的假想直线，
地球绕它自转，如同陀螺的转轴。
\emph{轨道}\index{轨道}是一个天体绕另一天体运行的闭合路径，
地球绕太阳、\omterm{def:g2:moon-in-the-sky:moon}{月球}绕地球的轨道——
都非常接近圆形。
\end{definition}

\begin{proposition}[三种运动]\label{prop:g6:sun-earth-moon:motions}
整个历法依赖于三种稳定的运动：
\begin{enumerate}
  \item 地球绕\omterm{def:g6:sun-earth-moon:axisorbit}{地轴}\emph{自转}一周，耗时 $24$
        \omterm{ex:g2:measuring-time:clock}{小时}——形成昼夜，让太阳与群星依次掠过天空；
  \item 地球绕太阳\emph{公转}一周，约需 $365$
        天又四分之一——即一年，使太阳的弧线
        上下起伏，形成四季之变；
  \item \omterm{def:g2:moon-in-the-sky:moon}{月球}绕地球\emph{公转}一周，约需一个
        月——其被阳光照亮的一半轮流朝向我们，
        使月相依次出现，正如你熟悉的月历所记。
\end{enumerate}
三种运动，三只时钟：日、月、年——
人类最古老的计时器，至今仍在运转。
\end{proposition}

\begin{omfigure}
\begin{tikzpicture}[scale=0.85]
  \fill[omThm] (5.6,2.2) circle (0.55);
  \foreach \a in {0,30,...,330} {
    \draw[omThm] (5.6,2.2) ++(\a:0.72) -- ++(\a:0.24);
  }
  \node[below, font=\small] at (5.6,1.4) {太阳};
  \draw[dashed, omExo, thick] (5.6,2.2) ellipse (4.6 and 1.75);
  \fill[omDef] (10.2,2.2) circle (0.24);
  \draw[-Stealth, thick] (10.2,2.75) arc (90:210:0.5);
  \node[above, font=\small] at (9.7,3.2) {自转：$24$ h};
  \draw[dashed, omExo] (10.2,2.2) circle (0.85);
  \fill[omExo] (10.8,2.8) circle (0.1);
  \node[right, font=\small] at (10.98,2.8) {月球：一个月};
  \draw[-Stealth, thick, omDef] (8.1,0.62) arc (222:262:3.9);
  \node[below, font=\small, omDef] at (8.3,0.35) {轨道：一年};
\end{tikzpicture}

{\small 系统示意（距离照例压缩）：地球
一边自转，一边绕太阳公转；\omterm{def:g2:moon-in-the-sky:moon}{月球}绕地球公转。}
\end{omfigure}

\begin{example}[四分之一天与闰年]\label{ex:g6:sun-earth-moon:leap}
地球绕行一周需要 $365$ 天\emph{又约四分之一天}——
大自然并不迁就整数。历法通过积攒这些零头来还账：
每四年积攒的四分之一天合成整整一天，
每隔四年加入历法一次——即 $366$ 天的
\emph{闰年}，多出的一天放在二月末。若忽略这笔账，
长达一个世纪，历法就会与\omterm{def:g3:sun-days-seasons:year}{季节}相差数周，
古老的历法为此付出了惨痛的代价。
\end{example}

\begin{example}[按比例的系统]\label{ex:g6:sun-earth-moon:scale}
笔记本上的数字：地球直径约 \qty{12800}{km}，
\omterm{def:g2:moon-in-the-sky:moon}{月球}直径约为其四分之一，绕地球运行，距地球约
\qty{384000}{km}，约等于三十个地球首尾相接的长度。太阳
直径约为 $109$ 个地球，距我们 $150$ 百万
千米。把地球缩成直径一\omterm{def:g2:measuring-length:units}{厘米}的弹珠，模型显示：
\omterm{def:g2:moon-in-the-sky:moon}{月球}是一粒胡椒，相距 $30$ \omterm{def:g2:measuring-length:units}{厘米}；太阳是一只
直径超过一米的大球，立在街道另一头一百多米之外。
我们身边的天地大部分是空的——正如公园里的
\omterm{def:g4:solar-system:family}{太阳系}比例模型早已悄悄暗示过这一点。
\end{example}

\section{还清的欠账：四季的成因}

\begin{proposition}[倾斜造成四季]\label{prop:g6:sun-earth-moon:tilt}
地球的\omterm{def:g6:sun-earth-moon:axisorbit}{地轴}在\omterm{def:g6:sun-earth-moon:axisorbit}{轨道}上并不竖直：它倾斜约
直角的四分之一——并且在整个公转过程中始终
\emph{朝同一方向}倾斜，好似天空中一只忠实的\omterm{def:g4:magnets-and-compass:compass}{指南针}。
于是，在公转的一半时间里，我们所在的半球
\emph{朝向}太阳倾斜：太阳高悬，白昼变长——
这便是夏季。半圈之后，倾斜方向不变，
但我们所在的半球\emph{背离}太阳：太阳低垂，白昼变短——冬季。
四季并非因太阳变化，也非因距离远近：而是一颗倾斜的\omterm{def:g4:solar-system:planet}{行星}，
带着不变的倾斜绕太阳运转。
\end{proposition}

\begin{omfigure}
\begin{tikzpicture}[scale=0.85]
  \fill[omThm] (5.6,2.0) circle (0.5);
  \foreach \a in {0,45,...,315} {
    \draw[omThm] (5.6,2.0) ++(\a:0.65) -- ++(\a:0.22);
  }
  \begin{scope}[shift={(1.1,2.0)}]
    \draw[thick] (0,0) circle (0.62);
    \draw[very thick, omDef] (-0.28,-0.75) -- (0.28,0.75);
    \fill[omDef] (0.19,0.5) circle (0.09);
  \end{scope}
  \node[below, font=\small, align=center] at (1.1,0.9)
    {倾向太阳：\\此处为夏季};
  \begin{scope}[shift={(10.1,2.0)}]
    \draw[thick] (0,0) circle (0.62);
    \draw[very thick, omDef] (-0.28,-0.75) -- (0.28,0.75);
    \fill[omDef] (0.19,0.5) circle (0.09);
  \end{scope}
  \node[below, font=\small, align=center] at (10.1,0.9)
    {半圈之后，倾斜依旧：\\背离太阳——此处为冬季};
\end{tikzpicture}

{\small 倾斜的地球位于其公转\omterm{def:g6:sun-earth-moon:axisorbit}{轨道}上的两个位置。\omterm{def:g6:sun-earth-moon:axisorbit}{地轴}
（蓝色）在两处都指向同一方向：六月时朝向太阳，
十二月时背离太阳，此处为标出的北方城镇。}
\end{omfigure}

\begin{method}[地球仪与灯]\label{met:g6:sun-earth-moon:globe}
在桌上看整个装置：
\begin{enumerate}
  \item 将一盏灯放在桌子中央作为太阳，手持倾斜的地球仪
        （或斜插在竹签上的球）绕它缓慢转动，
        整个过程中\emph{始终让竹签指向}
        房间里的同一个角落；
  \item 每转过公转的四分之一圈，停下来，看看你家所在的
        纬度：看每天的自转圈有多大部分处在
        光照之中；
  \item 灯的一侧，你所在的半球倾向灯光，
        ——光照弧长，灯位高：夏季；另一侧倾向背离，
        ——光照弧短，灯位低：冬季；
  \item 处于两者之间的两个位置，两极都不倾向灯光：昼夜
        平分——春季和秋季。
\end{enumerate}
让它奏效的唯一要求：决不让竹签摆动
——\omterm{def:g6:sun-earth-moon:axisorbit}{地轴}的指向始终不变，\emph{正是}\omterm{def:g3:sun-days-seasons:year}{季节}的成因。
\end{method}

\begin{example}[旧线索的解释]\label{ex:g6:sun-earth-moon:clues}
你多年来用\omterm{def:g3:sun-days-seasons:sundial}{日晷}观察到的每一个现象，现在都可归因于\omterm{def:g6:sun-earth-moon:axisorbit}{地轴}的倾斜。
夏季的高弧线：我们所在的半球倾向太阳，因此太阳
正午位置更高。冬季短粗的弧线和长长的\omterm{def:g1:light-and-shadows:shadow}{影子}：背离太阳，
太阳掠过低空。跨过赤道的两个半球\omterm{def:g3:sun-days-seasons:year}{季节}相反——
十二月南半球度假看海滩，北半球却在铲雪：
一个半球倾向太阳时，另一个半球必然背离太阳。就连
春秋两季昼夜等长，也可从灯与地球仪的
演示中推出。一个四分之一直角的倾斜，支配着整个过程。
\end{example}

\begin{remark}[不是距离——以及证明]\label{rem:g6:sun-earth-moon:notdistance}
地球\omterm{def:g6:sun-earth-moon:axisorbit}{轨道}非常接近圆形，因此朝向太阳的距离
变化很小——有趣的是，我们在北半球冬季时
反而略微\emph{更靠近}太阳。\omterm{def:g3:sun-days-seasons:year}{季节}关心的是
\omterm{def:g6:light-propagation:ray}{光线}的\emph{倾斜}，而不是路程的长短：
倾向太阳时，阳光以陡峭而集中的方式照射；背离太阳时，
同样的光则稀薄地铺在地面上。把手电筒
分别正对墙面和斜照墙面，一秒钟内就能看出差别——
陡射的光灼热，掠射的光几乎不暖。
\end{remark}

\section{练习}

\begin{exercise}[$\star$]\label{exo:g6:sun-earth-moon:1}
说出该系统的三种运动，并指出每一种运动各“拨动”
哪一个时钟。
\end{exercise}

\begin{exercise}[$\star$]\label{exo:g6:sun-earth-moon:2}
什么是\omterm{def:g6:sun-earth-moon:axisorbit}{地轴}？关于\omterm{def:g6:sun-earth-moon:axisorbit}{地轴}，哪一个“四分之一直角”的事实，
加上怎样的指向恒定性，才造成\omterm{def:g3:sun-days-seasons:year}{季节}？
\end{exercise}

\begin{exercise}[$\star$]\label{exo:g6:sun-earth-moon:3}
为什么历法每四年要增加一天？如果从不增加，
会出现什么偏差？
\end{exercise}

\begin{exercise}[$\star$]\label{exo:g6:sun-earth-moon:4}
大约要把多少个地球排成一列，才能到达\omterm{def:g2:moon-in-the-sky:moon}{月球}？
又需要多少个地球并排，才能跨越太阳的圆面？
\end{exercise}

\begin{exercise}[$\star$]\label{exo:g6:sun-earth-moon:5}
在地球仪与灯的演示中，为什么竹签必须始终指向
房间里的同一个角落？若它摆动，会破坏什么？
\end{exercise}

\begin{exercise}[$\star$]\label{exo:g6:sun-earth-moon:6}
当此处为夏季时，赤道另一侧的半球是什么\omterm{def:g3:sun-days-seasons:year}{季节}，
为什么必然如此？
\end{exercise}

\begin{exercise}[$\star$]\label{exo:g6:sun-earth-moon:7}
利用手电筒照射墙面的图示（见
\cref{rem:g6:sun-earth-moon:notdistance}），解释为何倾向太阳时
我们获得的\omterm{def:g5:heat-and-insulation:heat}{热量}更多。
\end{exercise}

\begin{exercise}[$\star$]\label{exo:g6:sun-earth-moon:8}
一位老朋友又来了：为什么\omterm{def:g2:moon-in-the-sky:moon}{月球}在一个
月中会呈现月相变化？三种运动中，哪一种运动主导了这一变化？
\end{exercise}

\begin{exercise}[$\star\star$]\label{exo:g6:sun-earth-moon:9}
一位固执的叔叔坚称：“冬季是地球离太阳
最远的时候。”请从本章中列出两条共同
可以驳倒他的事实。
\end{exercise}

\begin{exercise}[$\star\star$]\label{exo:g6:sun-earth-moon:10}
假设\omterm{def:g6:sun-earth-moon:axisorbit}{地轴}在\omterm{def:g6:sun-earth-moon:axisorbit}{轨道}上完全竖直站立——
毫无倾斜。描述一年的情况：\omterm{def:g3:sun-days-seasons:year}{季节}、正午太阳
高度和白昼长短会发生什么变化？（在脑中让竖直的
地球仪绕灯走一圈。）
\end{exercise}

\begin{exercise}[$\star\star$]\label{exo:g6:sun-earth-moon:11}
\omterm{def:g2:moon-in-the-sky:moon}{月球}绕行一圈约需 $27$ 天，而从\omterm{ex:g2:moon-in-the-sky:shapes}{满月}到\omterm{ex:g2:moon-in-the-sky:shapes}{满月}
约需 $29.5$ 天。请提出解释——\omterm{def:g2:moon-in-the-sky:moon}{月球}绕行时，
还有什么也在运动？（画出在\omterm{def:g2:moon-in-the-sky:moon}{月球}绕行一圈期间
地球沿其\omterm{def:g6:sun-earth-moon:axisorbit}{轨道}前进的示意图。）
\end{exercise}

\begin{exercise}[$\star\star\star$]\label{exo:g6:sun-earth-moon:12}
仲夏时节，\omterm{def:g4:magnets-and-compass:poles}{北极}附近的地区在午夜时
太阳仍高悬于地平线之上——即著名的“午夜太阳”——而在仲冬时节，
那里的人们要等上数周才能迎来\omterm{def:g1:day-and-night:daynight}{日出}。请借助
倾斜的自转圆周（见
\cref{met:g6:sun-earth-moon:globe}）解释这两种奇观：六月与十二月，
倾斜对\omterm{def:g4:magnets-and-compass:poles}{北极}附近城镇每日圆周运动分别有何影响？
\end{exercise}

\section{问题：操场模型}

\begin{problem}[{周末问题——全班在操场上搭建
太阳–地球–月球系统；弹珠、胡椒粒
与一段很长的步行；倾斜也纳入模型}]\label{pb:g6:sun-earth-moon:1}
全班按比例搭建这一系统：地球是一颗
直径为 \qty{1}{cm} 的弹珠。备用的真实数据：地球直径 \qty{12800}{km}；
\omterm{def:g2:moon-in-the-sky:moon}{月球}直径 \qty{3200}{km}，绕地运行的距离为 \qty{384000}{km}；
太阳直径是地球的 $109$ 倍，距地球 $150$ 百万千米。

\medskip\noindent\textbf{第一部分——准备部件。}
\begin{enumerate}
  \item 在模型中，\qty{1}{cm} 代表多少千米？
  \item \omterm{def:g2:moon-in-the-sky:moon}{月球}的宽度约为地球的四分之一。模型中的
        \omterm{def:g2:moon-in-the-sky:moon}{月球}有多宽？哪种厨房用品
        最适合充当它？
  \item \omterm{def:g2:moon-in-the-sky:moon}{月球}的真实距离为地球直径的 $30$ 倍。模型中的
        \omterm{def:g2:moon-in-the-sky:moon}{月球}应距弹珠多远？
  \item 模型中的太阳直径为 $109$ 颗弹珠。这有多宽，
        \omterm{def:g6:measurement-in-science:measuring}{单位}用\omterm{def:g2:measuring-length:units}{厘米}表示？哪种操场器材的大小
        与之相近？
\end{enumerate}

\medskip\noindent\textbf{第二部分——漫长的步行。}
\begin{enumerate}[resume]
  \item 太阳距地球 $150$ 百万千米。在模型中——
        用第 1 题的答案去除，
        这段距离是多少\omterm{def:g2:measuring-length:units}{厘米}？合多少米？
  \item 模型能放进一块 \qty{100}{m} 的操场吗？
        太阳球应放在何处？
  \item 一名学生以 $1$ 米每秒的\omterm{def:g5:speed-distance-time:speed}{速度}走完
        地球–太阳这段距离，大约需要多少分钟？
  \item 站在弹珠旁，回望距离 $117$ 米的一米宽
        太阳球，全班颇感惊讶：它看起来竟与距弹珠
        $30$ \omterm{def:g2:measuring-length:units}{厘米}处、直径四分之一\omterm{def:g2:measuring-length:units}{厘米}的\omterm{def:g2:moon-in-the-sky:moon}{月球}颗粒
        大小相仿。请核对这两个比值，
        并把这一巧合与真实天空中偶尔出现的一种罕见而
        壮观的天象联系起来。
        以真实天空中的方式，壮观而罕见地出现。
\end{enumerate}

\medskip\noindent\textbf{第三部分——让模型动起来。}
\begin{enumerate}[resume]
  \item 若要如实地让模型动起来，
        弹珠绕太阳球转一圈要多长时间，豌豆绕弹珠
        转一圈要多长时间（现实中一天按一秒计）？
        （粗略估算时，一年按 $365$ 天、\omterm{def:g2:moon-in-the-sky:moon}{月球}绕行一圈按 $27$
        天计算。）
  \item 弹珠还必须自转。若现实中一天对应一秒，
        模型中它转得有多快？为什么在
        \qty{1}{cm}的弹珠上没有人能看出它在转动？
  \item 穿过弹珠的竹签充当\omterm{def:g6:sun-earth-moon:axisorbit}{地轴}。竹签应以多大倾角，
        并朝向何方，
        才能在整个绕行过程中被正确携带？
  \item 写出模型的诚实声明：指出一件
        操场模型能真实展示的事物，以及一件
        任何操场大小的模型都无法同时展示的内容（大小与距离
        能否同时？运动？\omterm{def:g6:light-propagation:ray}{光线}？）。请用你算出的数据加以说明。
\end{enumerate}
\end{problem}
